% ============================================================
% FOREWORD
% front_matter/foreword.tex
% ============================================================

\chapter*{Foreword}
\addcontentsline{toc}{chapter}{Foreword}
\markboth{Foreword}{Foreword}

Physics is the heart of all the known sciences. Its basic understanding is essential to both the natural and the developmental sciences, and every student who aspires to understand nature and its workings needs a thorough grasp of the subject to succeed in their chosen field.

In this regard, ``The Physics Codex: A Complete University Foundation'' approaches the subject from first principles. It rests on an important educational idea: a strong foundation in physics comes from understanding the reasoning behind the equations, not from memorizing them.

The book's greatest strength is how it explains things. It builds intuition before introducing mathematical expressions, so students learn to understand physical principles rather than recite formulas. Its derivations, diagrams, worked examples, exercises, and discussions of common misconceptions give readers a structured and engaging learning experience.

The book covers five major areas of physics: Mechanics; Fluids and Thermal Physics; Waves and Optics; Electricity and Magnetism; and Modern Physics. Its broad scope, together with its focus on bridging secondary-school knowledge and first-year university requirements, makes it a valuable resource for students preparing for advanced scientific study.

The book is also relevant to science education in Nigeria. Accessible explanations and strong conceptual foundations can help students approach physics with greater confidence. Its real-world applications connect scientific principles to familiar experiences without sacrificing intellectual rigour.

I believe ``The Physics Codex'' will benefit Senior Secondary School students preparing for the Unified Tertiary Matriculation Examination (UTME), university students strengthening their foundations, and teachers seeking additional instructional resources. I commend Korede Samuel Omotosho for this educational initiative, and I hope the book inspires its readers to approach physics with curiosity, discipline, and a deeper appreciation of the natural world.

\vspace{2cm}

\begin{flushright}
\textit{Engr. Edidi S. I. Ayodeji}\\
\textit{B.Eng., M.Sc., MNSE, COREN}\\
\textit{MD, Envostructs Nigeria Limited}\\
\textit{(Former HOD, Civil Engineering, Lagos State Polytechnic, Ikorodu)}\\
\textit{[Lagos, Nigeria]}\\
\textit{[October, 2026]}
\end{flushright}